\section{Reproducibility, provenance, and integration artifacts}
\label{sec:reproducibility}

\subsection{What the archive contains}

The archive is self-contained for the Python checks and the manuscript
build. Its top-level structure is shown below; the included provenance
and hash manifests identify the exact source snapshot and additions.
The scanner requires Python 3.10 or later and has no third-party runtime
dependencies. The article uses a conventional pdf\LaTeX/Bib\TeX build.
Matplotlib is needed only to regenerate the optional static figure, whose
rendered PDF is already included.

\begin{center}
\small
\begin{tabularx}{\textwidth}{@{}p{0.30\textwidth}X@{}}
\toprule
Path & Contents and purpose \\
\midrule
\path{article/} & This PDF, modular \LaTeX{} sources, bibliography,
figure source, and build script. \\
\path{fast/} & The pinned scanner snapshot plus the new graded,
corridor, and scalar-component modules, public integration, tests,
and benchmark drivers. \\
\path{dihedral_covers/} & Standalone arithmetic covering kernel,
independent expanded-sheet tests, examples, and its mathematical
input/output contract. \\
\path{reference/} & Frozen preceding forward-transfer prototype and
the initial corridor implementation used by the primary timing run. \\
\path{verification/} & Independent contraction audit, recorded results,
full test logs, and experiment interpretations. \\
\path{integration/} & Patch against the pinned repository baseline
and review instructions. \\
\path{verify.py} & Runs the production suite and independent audits in
a temporary copy, preserving the archived measurements. \\
\path{reproduce_benchmarks.py} & Replays either timing experiment in
an isolated copy with the intended implementation version. \\
\bottomrule
\end{tabularx}
\end{center}

The complete scanner copy makes the package runnable before repository
integration. The patch makes the precise production changes separately
reviewable. Prior research reports remain authoritative for inherited
features; this package does not silently replace their proofs or
renumber the repository's report series.

\subsection{Running the checks}

From the unpacked package root:
\begin{lstlisting}
python verify.py
\end{lstlisting}
The script first verifies the archived file hashes, then runs the scanner
tests, the dihedral tests, the independent contraction comparison, the
graded-support audit, and the scalar-representation audit. Temporary
results are compared with the deterministic counts in the recorded
reports. Elapsed times are intentionally not required to match.

The principal recorded correctness checks are:
\begin{itemize}
\item The complete scanner suite: 268 tests passed, including the 245
baseline tests and 23 new tests for grading, corridor transfer, scalar
components, and integration behavior.
\item The covering kernel: 10 test groups passed, including 13,846
expanded-sheet topology comparisons. Separate checks test conjugation,
explicit generic-cover equivariance, exhaustive exceptional-fibre
bijections, degenerate one- and two-sheet actions, and binary sheet
counts as large as $2^{10000}$.
\item The independent contraction audit: 80 actual diagrams, 555 prefix
stages, and 3,885 entrywise comparisons across seven configurations.
Every stage checks the eight special-contraction identities of the
frozen reference; 185 stages have already undergone partial scalar
cancellation, and 377 have nonzero minimal differentials.
\item The support audit: 72 diagrams and 639 stages; both the existing
degree-gap certificate and the stronger full support certificate
accept the same 196 stages in this experiment.
\item The scalar-representation audit: exact equality of $i,p,h$ with
the frozen reference on five shared-suffix sizes, together with the
proved packed-bit and sparse-support formulas.
\end{itemize}

The independent oracles are deliberately different from the optimized
representations. Cover topology is obtained by enumerating small sheet
actions, tracing boundary cycles, and propagating orientation parity.
The transfer comparison uses the preceding complete forward formula
and verifies its actual coefficient matrices. Additional scanner tests
compare with the separate set-coefficient, LIFO-pivot backend. These
checks are strong finite evidence of implementation correctness; the
general claims rest on the written proofs.

\subsection{Replaying the timing experiments}

The two recorded timing datasets have different scopes and should not be
pooled. To rerun them without altering archived source or result files:
\begin{lstlisting}
python reproduce_benchmarks.py --experiment baseline \
  --output baseline_rerun.json
python reproduce_benchmarks.py --experiment compressed \
  --output compressed_rerun.json
\end{lstlisting}

The baseline replay restores the frozen first corridor version in a
temporary copy before invoking the original nine-arm benchmark. The
compressed replay uses the current source and the frozen version as
its explicit comparison arm. Both default to seven shuffled paired
rounds and a 20 ms pilot batch target. Original raw samples, repetitions,
operation counters, Python/platform information, random seeds, and
implementation hashes are included in the recorded JSON files.

Run the benchmarks with other CPU-intensive work paused. Timing ratios
are observations on the recorded environment, not portable complexity
theorems. The unchanged standard/control arms provide an internal
measure of drift; the exact planted-output and operation-count checks
should agree even when measured times do not. The benchmark target is
the full scanner for actual diagrams and the complete reduction stage
for the explicitly constructed algebraic examples. The latter does not
include creation of the synthetic input.

\subsection{Building the paper and reviewing the patch}

The manuscript can be rebuilt with:
\begin{lstlisting}
cd article
sh build.sh
\end{lstlisting}
The build requires pdf\LaTeX{}, Bib\TeX{}, and the standard packages
listed in the preamble. The bibliography supplies primary literature
and explicit provenance for the earlier local prototype. The
manuscript was compiled to PDF, inspected through rendered pages,
and checked for unresolved references and content overflow before
packaging.

At the pinned repository commit, inspect the supplied patch and run
\begin{lstlisting}
git apply --check /path/to/scanner_changes.patch
\end{lstlisting}
from the repository root. The integration guide explains how to retain
the frozen test reference next to the scanner copy or adapt its test
location consistently. The article and standalone geometric kernel
can then be placed in the repository's chosen research-report
directory. No report number is assumed, and no remote commit or
publication is performed by this package.

\subsection{Limits of the evidence}

The reported tests do not amount to exhaustive verification of all
knot diagrams or a proof-assistant formalization. The synthetic
separation families are valid graded complexes in the implemented
coefficient category; their realization as knot prefixes is a stated
open question. The ordinary fixture timings do not support replacing
the default sparse reducer. The geometric kernel requires its explicit
covering promise and external attachment data for use in a hierarchy.
These qualifications identify the boundaries of the delivered results
and the concrete work still required for a global complexity claim.
